\subsubsection{\textit{Score} Geral de Engajamento}
O Score Geral de Engajamento (\textit{Engagement Score}) é uma métrica composta que integra múltiplas dimensões de participação em um único valor numérico entre 0 e 1. Esta medida sintética permite uma avaliação holística do engajamento estudantil, combinando indicadores de foco, presença temporal e participação ativa através de uma abordagem ponderada.

O cálculo é realizado pela função \textbf{calculate\_engagement\_score(metrics\_dict)} que implementa uma média ponderada de cinco componentes fundamentais:

\begin{equation}
\text{Engagement Score} = \sum_{i=1}^{5} w_i \times m_i
\end{equation}

\textbf{onde:}
\begin{itemize}
    \item $w_i$ é o peso atribuído ao componente $i$
    \item $m_i$ é o valor normalizado da métrica componente $i$
\end{itemize}

% \lstinputlisting[
%     language=python,
%     caption={Cálculo do \textit{Score} de Engajamento},
%     firstline=1,
%     lastline=46,
%     label={lst:calculate_engagement_score}
% ]{scripts/metricsCalc/scores.py}

\textbf{Componentes do score e seus pesos}:
\begin{table}[ht]
    \centering
    \vspace{0.5cm}
    \begin{tabular}{|c|c|l|}
        \hline
        Componente & Peso & Descrição \\
        \hline
        Foco & 35\% & Tempo na aba da aula (\texttt{lecture\_focus\_ratio}) \\
        % \hline
        Presença & 30\% & Razão tempo médio/duração total (\texttt{attendance\_ratio}) \\
        % \hline
        Participação & 20\% & Sinais ativos de engajamento (\texttt{voluntary\_participation}) \\
        % \hline
        Câmera & 10\% & Tempo com câmera ativa (\texttt{camera\_engagement}) \\
        % \hline
        Microfone & 5\% & Tempo com microfone ativo (\texttt{mic\_engagement}) \\
        \hline
    \end{tabular}
    \caption{Componentes e pesos do Score Geral de Engajamento}
    \label{tab:engagement_score_components}
\end{table}

\textbf{Justificativa dos pesos}:
\begin{itemize}
    \item \textbf{Foco na Aula (35\%)}: Considerado o indicador mais direto de atenção cognitiva
    \item \textbf{Presença Temporal (30\%)}: Base fundamental para qualquer tipo de engajamento
    \item \textbf{Participação Voluntária (20\%)}: Indica iniciativa ativa do estudante
    \item \textbf{Engajamento por Câmera (10\%)}: Contribuição visual limitada por fatores contextuais
    \item \textbf{Engajamento por Microfone (5\%)}: Participação vocal menos frequente em aulas expositivas
\end{itemize}


\subsubsection{Saúde da Atenção}
A Saúde da Atenção (\textit{Attention Health}) é uma métrica composta que avalia a qualidade do padrão de foco dos estudantes, considerando não apenas a quantidade de tempo focado, mas também características qualitativas como sustentação, consistência e baixa fragmentação. Este score, variando de 0 a 1, mede a robustez do processo atencional durante as sessões de aprendizagem.

O cálculo é realizado pela função \textbf{calculate\_attention\_health(metrics\_dict)} que integra quatro dimensões qualitativas da atenção:

\begin{equation}
\text{Attention Health} = \sum_{i=1}^{4} w_i \times h_i
\end{equation}

\textbf{onde:}
\begin{itemize}
    \item $w_i$ é o peso atribuído à dimensão $i$
    \item $h_i$ é o valor normalizado da dimensão de saúde $i$
\end{itemize}


% \lstinputlisting[
%     language=python,
%     caption={Cálculo da Saúde da Atenção},
%     firstline=49,
%     lastline=110,
%     label={lst:calculate_attention_health}
% ]{scripts/metricsCalc/scores.py}

\textbf{Dimensões da saúde da atenção e seus pesos}:
\begin{table}[H]
    \centering
    \vspace{0.5cm}
    \begin{tabular}{|c|c|l|}
        \hline
        Dimensão & Peso & Descrição \\
        \hline
        Duração do Foco & 30\% & Capacidade de manter foco sustentado \\
        Baixa Fragmentação & 25\% & Minimização de interrupções no foco \\
        Controle de Multitarefa & 25\% & Gestão adequada de múltiplas tarefas\\
        Consistência Temporal & 20\% & Estabilidade do engajamento ao longo do tempo\\
        \hline
    \end{tabular}
    \caption{Dimensões e pesos da Saúde da Atenção}
    \label{tab:attention_health_components}
\end{table}

\textbf{Funções auxiliares}:
\begin{itemize}
    \item \textbf{calculate\_focus\_duration\_health()}: Avalia a adequação da duração dos períodos de foco (nem muito curtos, nem muito longos)
    \item \textbf{calculate\_fragmentation\_health()}: Mede o grau de continuidade atencional (menor fragmentação = mais saudável)
    \item \textbf{calculate\_multitasking\_health()}: Avalia a gestão de múltiplas tarefas (multitasking produtivo vs. destrutivo)
    \item \textbf{calculate\_consistency\_health()}: Analisa a estabilidade do engajamento ao longo da sessão
\end{itemize}

\subsubsection{Risco de Distração}
O Risco de Distração (\textit{Distraction Risk}) é uma métrica composta que quantifica a propensão dos estudantes à desatenção durante as sessões de aprendizagem. Este score, variando de 0 a 1, integra múltiplos indicadores comportamentais para fornecer uma medida sintética do nível de distração, sendo valores mais altos indicativos de maior propensão à perda de foco.

O cálculo é realizado pela função \textbf{calculate\_distraction\_risk(metrics\_dict)} que implementa uma média ponderada de quatro componentes críticos:

\begin{equation}
\text{Distraction Risk} = \sum_{i=1}^{4} w_i \times f_i
\end{equation}

\textbf{onde:}
\begin{itemize}
    \item $w_i$ é o peso atribuído ao fator de risco $i$
    \item $f_i$ é o valor normalizado do fator de risco $i$
\end{itemize}


% \lstinputlisting[
%     language=python,
%     caption={Cálculo do Risco de Distração},
%     firstline=170,
%     lastline=210,
%     label={lst:calculate_distraction_risk}
% ]{scripts/metricsCalc/scores.py}

\textbf{Fatores de risco e seus pesos}:
\begin{table}[ht]
    \centering
    \vspace{0.5cm}
    \begin{tabular}{|c|c|l|}
        \hline
        Fator de Risco & Peso & Descrição \\
        \hline
        Proporção de Distração & 40\% & Percentual de tempo em sites distractores\\
        Frequência de Distração & 25\% & Número de transições para distração por hora\\
        Força das Distrações & 20\% & Impacto dos distractores principais\\
        Intensidade de Trocas & 15\% & Ritmo de alternância entre abas\\
        \hline
    \end{tabular}
    \caption{Fatores de risco e pesos do Score de Distração}
    \label{tab:distraction_risk_components}
\end{table}

\textbf{Funções auxiliares}:
\begin{itemize}
    \item \textbf{calculate\_main\_distraction\_strength()}: Calcula o impacto relativo das principais categorias de distração detectadas
    \item \textbf{calculate\_tab\_switch\_intensity()}: Normaliza e pondera a frequência de troca de abas
\end{itemize}